% GENERATED FILE. Edit canonical lessons, then run npm run generate:books.
% canonical-source-hash: fnv1a64:915d5167a6287065
% canonical-lessons: SA-C313-upadravah, SA-C313-kalahah, SA-C313-ksobhah, SA-C313-vestayati, SA-C313-samayati

\chapter{Trouble in the Town}
\label{ch:sa-trouble-in-the-town}
\hlchaptermodality{313}

\begin{chapteropening}
\textbf{By the end of this chapter:} \emph{I can tell how trouble broke out and was calmed.}

Complete the last lesson of chapter 313: I can tell how trouble broke out and was calmed.
\end{chapteropening}

\section[\texorpdfstring{\sk{उपद्रवः} (upadravaḥ)}{उपद्रवः (upadravaḥ)}]{\sk{उपद्रवः} (upadravaḥ) — a disturbance; a riot}

\label{lesson:SA-C313-upadravah}



\begin{quote}
Before the new one: say the Sanskrit for a lie, then the Sanskrit for doubtful.
\end{quote}

\subsection*{You'll want to know: \sk{उपद्रवः}}
\textbf{\sk{उपद्रवः}} — \emph{upadravaḥ} — \textquotedblleft{}a disturbance; a riot\textquotedblright{}.

\begin{quote}
\sk{नगरे} \sk{सहसा} \sk{उपद्रवः} \sk{घटितः।}
\end{quote}

\begin{quote}
\emph{nagare sahasā upadravaḥ ghaṭitaḥ.}
\end{quote}

\begin{quote}
A disturbance suddenly broke out in the city.
\end{quote}

A masculine noun.

\begin{cousinweb}
\textbf{\sk{उपद्रवः}} (\emph{upadravaḥ}) is built from \emph{upa} and the root √dru, \textquotedblleft{}to run\textquotedblright{}; MW gives \emph{upadrava} as that which attacks suddenly, a misfortune, and also national commotion or rebellion (MW).
\end{cousinweb}

\begin{grammarlens}[title={What you've built}]
One more word for telling what happened.
\end{grammarlens}

\subsection*{Your turn}
\begin{itemize}
\raggedright
  \item \emph{Say it:} \emph{upadravaḥ}
  \item \emph{Say it:} \emph{upadravaḥ}, once more
  \item \emph{Say it:} \emph{nagare sahasā upadravaḥ ghaṭitaḥ.} — A disturbance suddenly broke out in the city.
  \item \emph{Recall:} say the Sanskrit for a lie, then the Sanskrit for doubtful, then say \emph{upadravaḥ} again
\end{itemize}

\begin{tcolorbox}[breakable,colback=teal!4,colframe=teal!35!black,title={Before you move on}]
What does \sk{उपद्रवः} mean? (\textquotedblleft{}A disturbance; a riot\textquotedblright{}.) Say it once more.
\end{tcolorbox}
\section[\texorpdfstring{\sk{कलहः} (kalahaḥ)}{कलहः (kalahaḥ)}]{\sk{कलहः} (kalahaḥ) — a quarrel}

\label{lesson:SA-C313-kalahah}



\begin{quote}
Before the new one: say the Sanskrit for doubtful, then the Sanskrit for a disturbance.
\end{quote}

\subsection*{You'll want to know: \sk{कलहः}}
\textbf{\sk{कलहः}} — \emph{kalahaḥ} — \textquotedblleft{}a quarrel\textquotedblright{}.

\begin{quote}
\sk{विपणौ} \sk{कलहः} \sk{भवति।}
\end{quote}

\begin{quote}
\emph{vipaṇau kalahaḥ bhavati.}
\end{quote}

\begin{quote}
There is a quarrel in the market.
\end{quote}

A masculine noun.

\begin{cousinweb}
\textbf{\sk{कलहः}} (\emph{kalahaḥ}) MW gives \emph{kalaha} as strife, contention, a quarrel or fight, without a derivation, so its origin is uncertain (MW).
\end{cousinweb}

\begin{grammarlens}[title={What you've built}]
One more word for telling what happened.
\end{grammarlens}

\subsection*{Your turn}
\begin{itemize}
\raggedright
  \item \emph{Say it:} \emph{kalahaḥ}
  \item \emph{Say it:} \emph{kalahaḥ}, once more
  \item \emph{Say it:} \emph{vipaṇau kalahaḥ bhavati.} — There is a quarrel in the market.
  \item \emph{Recall:} say the Sanskrit for doubtful, then the Sanskrit for a disturbance, then say \emph{kalahaḥ} again
\end{itemize}

\begin{tcolorbox}[breakable,colback=teal!4,colframe=teal!35!black,title={Before you move on}]
What does \sk{कलहः} mean? (\textquotedblleft{}A quarrel\textquotedblright{}.) Say it once more.
\end{tcolorbox}
\section[\texorpdfstring{\sk{क्षोभः} (kṣobhaḥ)}{क्षोभः (kṣobhaḥ)}]{\sk{क्षोभः} (kṣobhaḥ) — unrest; agitation}

\label{lesson:SA-C313-ksobhah}



\begin{quote}
Before the new one: say the Sanskrit for a disturbance, then the Sanskrit for a quarrel.
\end{quote}

\subsection*{You'll want to know: \sk{क्षोभः}}
\textbf{\sk{क्षोभः}} — \emph{kṣobhaḥ} — \textquotedblleft{}unrest; agitation\textquotedblright{}.

\begin{quote}
\sk{नगरे} \sk{क्षोभः} \sk{वर्धते।}
\end{quote}

\begin{quote}
\emph{nagare kṣobhaḥ vardhate.}
\end{quote}

\begin{quote}
Unrest is growing in the city.
\end{quote}

A masculine noun.

\begin{cousinweb}
\textbf{\sk{क्षोभः}} (\emph{kṣobhaḥ}) comes from the root √kṣubh, \textquotedblleft{}to shake, be agitated\textquotedblright{}; MW gives \emph{kṣobha} as shaking, agitation, disturbance (MW).
\end{cousinweb}

\begin{grammarlens}[title={What you've built}]
One more word for telling what happened.
\end{grammarlens}

\subsection*{Your turn}
\begin{itemize}
\raggedright
  \item \emph{Say it:} \emph{kṣobhaḥ}
  \item \emph{Say it:} \emph{kṣobhaḥ}, once more
  \item \emph{Say it:} \emph{nagare kṣobhaḥ vardhate.} — Unrest is growing in the city.
  \item \emph{Recall:} say the Sanskrit for a disturbance, then the Sanskrit for a quarrel, then say \emph{kṣobhaḥ} again
\end{itemize}

\begin{tcolorbox}[breakable,colback=teal!4,colframe=teal!35!black,title={Before you move on}]
What does \sk{क्षोभः} mean? (\textquotedblleft{}Unrest; agitation\textquotedblright{}.) Say it once more.
\end{tcolorbox}
\section[\texorpdfstring{\sk{वेष्टयति} (veṣṭayati)}{वेष्टयति (veṣṭayati)}]{\sk{वेष्टयति} (veṣṭayati) — surrounds}

\label{lesson:SA-C313-vestayati}



\begin{quote}
Before the new one: say the Sanskrit for a quarrel, then the Sanskrit for unrest.
\end{quote}

\subsection*{You'll want to know: \sk{वेष्टयति}}
\textbf{\sk{वेष्टयति}} — \emph{veṣṭayati} — \textquotedblleft{}surrounds\textquotedblright{}.

\begin{quote}
\sk{आरक्षकाः} \sk{विपणिं} \sk{वेष्टयन्ति।}
\end{quote}

\begin{quote}
\emph{ārakṣakāḥ vipaṇiṃ veṣṭayanti.}
\end{quote}

\begin{quote}
The police surround the market.
\end{quote}

Past (he …): \sk{सः} \sk{विपणिं} \sk{वेष्टितवान्}, \emph{saḥ vipaṇiṃ veṣṭitavān}.

\begin{cousinweb}
\textbf{\sk{वेष्टयति}} (\emph{veṣṭayati}) comes from the root √veṣṭ, \textquotedblleft{}to wind or twist round\textquotedblright{}; MW gives its causative \emph{veṣṭayati} as to wrap up, enclose, surround (MW).
\end{cousinweb}

\begin{grammarlens}[title={What you've built}]
One more word for telling what happened.
\end{grammarlens}

\subsection*{Your turn}
\begin{itemize}
\raggedright
  \item \emph{Say it:} \emph{veṣṭayati}
  \item \emph{Say it:} \emph{veṣṭayati}, once more
  \item \emph{Say it:} \emph{ārakṣakāḥ vipaṇiṃ veṣṭayanti.} — The police surround the market.
  \item \emph{Recall:} say the Sanskrit for a quarrel, then the Sanskrit for unrest, then say \emph{veṣṭayati} again
\end{itemize}

\begin{tcolorbox}[breakable,colback=teal!4,colframe=teal!35!black,title={Before you move on}]
What does \sk{वेष्टयति} mean? (\textquotedblleft{}Surrounds\textquotedblright{}.) Say it once more.
\end{tcolorbox}
\section[\texorpdfstring{\sk{शमयति} (śamayati)}{शमयति (śamayati)}]{\sk{शमयति} (śamayati) — calms; quells}

\label{lesson:SA-C313-samayati}



\begin{quote}
Before the new one: say the Sanskrit for unrest, then the Sanskrit for surrounds.
\end{quote}

\subsection*{You'll want to know: \sk{शमयति}}
\textbf{\sk{शमयति}} — \emph{śamayati} — \textquotedblleft{}calms; quells\textquotedblright{}.

\begin{quote}
\sk{आरक्षकाः} \sk{उपद्रवं} \sk{शमयन्ति।}
\end{quote}

\begin{quote}
\emph{ārakṣakāḥ upadravaṃ śamayanti.}
\end{quote}

\begin{quote}
The police quell the disturbance.
\end{quote}

Past (he …): \sk{सः} \sk{कलहं} \sk{शमितवान्}, \emph{saḥ kalahaṃ śamitavān}.

\begin{cousinweb}
\textbf{\sk{शमयति}} (\emph{śamayati}) comes from the root √śam, \textquotedblleft{}to be quiet\textquotedblright{}, as in \emph{śāmyati}, in its causative form: to make quiet, calm (MW).
\end{cousinweb}

\begin{grammarlens}[title={What you've built}]
One more word for telling what happened.
\end{grammarlens}

\subsection*{Your turn}
\begin{itemize}
\raggedright
  \item \emph{Say it:} \emph{śamayati}
  \item \emph{Say it:} \emph{śamayati}, once more
  \item \emph{Say it:} \emph{ārakṣakāḥ upadravaṃ śamayanti.} — The police quell the disturbance.
  \item \emph{Recall:} say the Sanskrit for unrest, then the Sanskrit for surrounds, then say \emph{śamayati} again
\end{itemize}

\begin{tcolorbox}[breakable,colback=teal!4,colframe=teal!35!black,title={Before you move on}]
What does \sk{शमयति} mean? (\textquotedblleft{}Calms; quells\textquotedblright{}.) Say it once more.
\end{tcolorbox}
